\section{Discussão dos Resultados}
\label{sec:discussao}

\subsection{Impacto dos Padrões de Entrada no Desempenho}
\label{sec:impacto_padroes}

A Tabela~\ref{tab:comp_padroes} compara os 5 padrões na matriz Q4 ($40.000^2$). Optou-se por Q4 (3,2 GB) em vez de Q6 (20 GB) porque em matrizes extremas a saturação do barramento de memória homogeneíza os tempos (~20 s em todos os padrões), mascarando o custo da CPU. Em Q4, a sensibilidade da CPU aos desvios condicionais permanece nítida.

\begin{table}[H]
\centering
\small
\caption{Comparação de tempo sequencial e speedup entre os padrões em Q4 ($40.000^2$).}
\label{tab:comp_padroes}
\begin{tabular}{@{}llccccc@{}}
\toprule
\textbf{Modo} & \textbf{Padrão} & \textbf{Score} & $\overline{T}_{seq}$ (s) & \textbf{Speedup OMP (16t)} & \textbf{Speedup MPI (16p)} \\ \midrule
\multirow{5}{*}{\texttt{--full-matrix}} 
 & \texttt{p1}     & 39.999 & 1{,}936\,s & 2{,}93$\times$  & \textbf{4{,}08$\times$} \\
 & \texttt{p2}     & 0      & 3{,}881\,s & 5{,}62$\times$  & \textbf{6{,}58$\times$} \\
 & \texttt{p3}     & 10.000 & 3{,}405\,s & 4{,}99$\times$  & \textbf{6{,}07$\times$} \\
 & \texttt{p4}     & 40.000 & 2{,}280\,s & 3{,}44$\times$  & \textbf{4{,}67$\times$} \\
 & \texttt{random} & 33.744 & 2{,}092\,s & 3{,}17$\times$  & \textbf{4{,}50$\times$} \\ \midrule
\multirow{5}{*}{\texttt{--mem-opt}} 
 & \texttt{p1}     & 39.999 & 1{,}131\,s & 11{,}51$\times$ & \textbf{11{,}66$\times$} \\
 & \texttt{p2}     & 0      & 2{,}723\,s & \textbf{13{,}60$\times$} & 10{,}35$\times$ \\
 & \texttt{p3}     & 10.000 & 2{,}334\,s & \textbf{13{,}31$\times$} & 10{,}52$\times$ \\
 & \texttt{p4}     & 40.000 & 1{,}330\,s & \textbf{11{,}71$\times$} & 11{,}20$\times$ \\
 & \texttt{random} & 33.744 & 1{,}069\,s & 11{,}35$\times$ & \textbf{12{,}29$\times$} \\ \bottomrule
\end{tabular}
\end{table}

\textbf{O padrão altera drasticamente o tempo sequencial}: o padrão \texttt{p2} (conjuntos disjuntos) demora mais do que o \textbf{dobro} do tempo de \texttt{p1}. Quando há coincidência de caracteres (\texttt{p1} e \texttt{p4}), o processador realiza apenas uma leitura e uma adição. No caso disjunto (\texttt{p2}), o desvio falha sistematicamente, forçando duas leituras de memória distintas e a execução da instrução condicional de máximo (\texttt{cmov}), aumentando os ciclos por instrução (CPI). Quanto ao \textbf{speedup}, \texttt{p2} e \texttt{p3} atingem acelerações relativas maiores porque o maior custo computacional do miolo de cada célula amortiza com maior folga o overhead de despacho e sincronização das threads.

\subsection{Comparação com a Lei de Amdahl e Cache Blocking L1d}

Na matriz completa, a alocação sequencial e a saturação de largura de banda de memória (\textit{memory-bound}) barram o speedup em torno de $4\times \text{ a } 5{,}18\times$. No modo \texttt{--mem-opt}, a contenção desaparece e o sistema opera predominantemente como \textit{compute-bound}, atingindo \textbf{12{,}30x no OpenMP} (na matriz Q4) e \textbf{16{,}61x no MPI} (na matriz gigantesca Q8). 

Além disso, o MPI em 1 e 2 processos atinge \textbf{super-linearidade} ($S = 1{,}17\times \text{ a } 1{,}38\times$ com $p=1$, sustentando $16{,}61\times$ com $p=16$ em Q8, com eficiência de $1{,}04$): enquanto o sequencial varre a linha inteira estourando a cache de dados L1d ($32$\,KiB), o MPI fatia o domínio em \textit{chunks} de $1024$ colunas ($4$\,KiB), mantendo os dados integralmente na cache L1d e reduzindo expressivamente os \textit{cache misses}.

\subsection{Escalabilidade Forte e Fraca}

O algoritmo apresenta \textbf{alta escalabilidade forte} em \texttt{--mem-opt} para entradas de grande porte ($N \ge 40.000$), onde o speedup escala quase linearmente de 2 até 16 núcleos. Para entradas pequenas ($N \le 1.000$), a sobrecarga de gerência de tarefas inviabiliza ganhos paralelos. A \textbf{escalabilidade fraca} é excelente: ao escalar a área $N \cdot M$ proporcionalmente ao número de núcleos $p$, o tempo total mantém-se estável com variação mínima.

\subsection{Comportamento do SMT e Variabilidade Experimental}

As execuções foram vinculadas exclusivamente aos 16 núcleos físicos da CPU (\texttt{OMP\_PLACES=cores}, \texttt{OMP\_PROC\_BIND=close}), sem utilizar as 32 threads lógicas do SMT. Como o kernel do LCS é dominado por operações aritméticas simples em inteiros e leituras sequenciais, duas threads no mesmo núcleo físico disputariam as mesmas unidades de execução e a mesma cache L1d de 32 KiB, degradando o rendimento por contenção de recursos.

A dispersão entre as 20 repetições de cada teste foi mínima: o coeficiente de variação ($CV = \sigma / \mu$) apresentou mediana de 0{,}25\% e percentil 90 de 2{,}69\%. Esse baixo desvio padrão demonstra estabilidade nas medições e indica que o isolamento por afinidade de núcleos evitou perturbações decorrentes de trocas de contexto pelo sistema operacional.

\subsection{Impacto Geométrico: Quadrados vs. Retângulos e Retângulos Invertidos}
\label{sec:impacto_geometria}

A Tabela~\ref{tab:comp_geometria} compara o desempenho entre matrizes de mesma área de cálculo ($1{,}60 \times 10^9$ células) sob diferentes proporções de aspecto.

\begin{table}[H]
\centering
\small
\caption{Comparação direta de \textit{speedup} (16 núcleos) entre formatos com mesma área de cálculo.}
\label{tab:comp_geometria}
\begin{tabular}{@{}llcccc@{}}
\toprule
\textbf{Caso} & \textbf{Geometria / Razão} & \textbf{Dimensões} & \textbf{Área ($M \times N$)} & \textbf{OpenMP (16t)} & \textbf{MPI (16p)} \\ \midrule
Q4      & Quadrada ($1:1$)          & $40.000 \times 40.000$ & $1{,}60 \times 10^9$ & \textbf{12{,}30$\times$} & 11{,}21$\times$ \\
R2      & Retangular ($4:1$, direta)& $80.000 \times 20.000$ & $1{,}60 \times 10^9$ & \textbf{12{,}96$\times$} & 8{,}91$\times$ \\
R2\_inv & Retangular ($1:4$, inv.)  & $20.000 \times 80.000$ & $1{,}60 \times 10^9$ & \textbf{12{,}29$\times$} & 8{,}82$\times$ \\ \midrule
R3      & Retangular ($8:1$, direta)& $80.000 \times 10.000$ & $0{,}80 \times 10^9$ & \textbf{12{,}05$\times$} & 6{,}10$\times$ \\
R3\_inv & Retangular ($1:8$, inv.)  & $10.000 \times 80.000$ & $0{,}80 \times 10^9$ & \textbf{11{,}28$\times$} & 5{,}99$\times$ \\ \bottomrule
\end{tabular}
\end{table}

\textbf{Matriz Quadrada vs. Retangular:} Em matrizes de mesma área computacional (Q4 vs. R2), o MPI tem seu \textit{speedup} reduzido de $11{,}21\times$ para $8{,}91\times$, e para $6{,}10\times$ em R3 (8:1), pois com poucas linhas por processo o pipeline esvazia rápido e a comunicação domina. No OpenMP, a decomposição 2D mantém o paralelismo estável ($12{,}30\times$ em Q4 vs. $12{,}96\times$ em R2).

\textbf{Retângulos Diretos vs. Invertidos (\texttt{\_inv}):} No pipeline MPI, o caso direto ($80k \times 20k$) possui 78 \textit{chunks} concorrentes contra apenas 19 no invertido ($20k \times 80k$), reduzindo a sobreposição de mensagens ($8{,}91\times$ vs. $8{,}82\times$ em R2, e $6{,}10\times$ vs. $5{,}99\times$ em R3). No OpenMP, linhas mais longas em C favorecem a localidade contígua e o \textit{prefetcher} ($12{,}96\times$ vs. $12{,}29\times$).
